\clearpage
\SourceLeaf{175}%
\pdfbookmark[0]{Göttingische gelehrta Anzeigen. 1825 April 11.}{notice1825april11}
\begin{center}
\rule{.82\linewidth}{.35pt}\par
.Göttingische gelehrta Anzeigen.\quad 1825 April 11.\par
\rule{.82\linewidth}{.35pt}
\end{center}
On 5 April Court Councillor \textsc{Gauss} presented a lecture to the Royal Society, entitled:
\begin{center}\itshape
Theoria Residuorum Biquadraticorum, Commentatio prima.
\end{center}

The theory of quadratic residues is known to form one of the most interesting parts of higher arithmetic. After repeated investigations it can now be regarded as complete and closed. Accounts of its history can be found in these pages for 1808 May 12 and September 19, and 1817 March 10. The last of these also gave some preliminary information about the investigations which the author of the present paper had been pursuing since the year 1805 into the related theory of cubic and biquadratic residues, equally fruitful and interesting but far more difficult. Although he already possessed the essential elements of this theory at that time, other work has so far prevented him from publishing anything about it. Only now has he been able to begin writing up part of these investigations. He has begun with the theory of biquadratic residues, which is more closely related to the theory of quadratic residues than is that of cubic residues. The present paper, however, is
\SourceLeaf{176}%
by no means intended to exhaust this exceptionally rich subject. The development of the \emph{general} theory, which requires a quite distinctive extension of the domain of higher arithmetic, is reserved for a future continuation. This first paper contains those investigations which could be set out completely without such an extension. Only some of the results can be selected for this review.

An integer $a$ is called a biquadratic residue of the integer $p$ if there are numbers of the form $x^4-a$ divisible by $p$; a biquadratic nonresidue if no numbers of that form can be divisible by $p$. Clearly all biquadratic residues of $p$ are also quadratic residues of the same number, and hence all quadratic nonresidues are also biquadratic nonresidues. But not all quadratic residues are also biquadratic residues. It suffices to restrict the investigations to the case where $p$ is a prime of the form $4n+1$ and $a$ is not divisible by $p$, since all other cases can easily be reduced to this one.

The investigations of this subject divide into two parts, according as $p$ or $a$ is regarded as given. The first is much less difficult than the second and, compared with it, is to be regarded as entirely elementary. The paper contains a complete account of everything essential to be said about it.

From the second part, by contrast, only a few special cases are treated here. These could be dealt with without too much preparation and can serve as preliminaries to the general theory to be given in the future. They are the cases where $a=-1$ and $a=\pm2$ are set. \Needspace{10\baselineskip}The first presents no difficulty at all: the \emph{Disquisitiones Arithmeticae} had already shown that $-1$ is a biquadratic residue of $p$ whenever $p$ has the form $8n+1$, but only a quadratic residue and a biquadratic nonresidue of $p$ when $p$ has the form $8n+5$. The case $a=\pm2$ is quite different. It has long been known that $+2$ and $-2$ are quadratic, and therefore also biquadratic, nonresidues of $p$ when $p$ has the form $8n+5$, and at least quadratic residues when $p$ has the form $8n+1$. It is also known that for this form of $p$, $+2$ and $-2$ are either both biquadratic residues or both biquadratic nonresidues:
\SourceLeaf{177}%
but deciding which of these two cases occurs is an investigation of a much higher order. The paper develops two different criteria for this purpose.

The first criterion is connected with the decomposition of the number $p$ into one square plus twice another, which, as is known (since, as already remarked, $p$ is assumed to be prime), is always possible and possible in only one way. If we set $p=gg+2hh$, then $\pm2$ is a biquadratic residue of $p$ when $g$ has the form $8n+1$ or $8n+7$, but a biquadratic nonresidue when $g$ has the form $8n+3$ or $8n+5$.

The second criterion is connected with the decomposition of the number $p$ into two squares, which is also known to be always possible and possible in only one way. If we set $p=ee+ff$ and suppose that $ee$ denotes the odd square and $ff$ the even square, then the assumed form $p=8n+1$ already implies that $\frac12f$, too, is even. Thus $f$ has either the form $8m$ or the form $8m+4$: in the first case $\pm2$ is a biquadratic residue, in the other a biquadratic nonresidue of $p$.

We only touch here on an observation for which higher arithmetic so often provides an occasion: it is not so much the beauty and simplicity of the theorems themselves as the difficulty of proving them that makes them especially remarkable. Once there is reason to suspect a connection between the behaviour of the number $\pm2$ and the two decompositions of the number $p$ just described, it is extremely easy to discover that connection by examining examples. But even for the first criterion its proof is not quite easy; for the second it lies much deeper, closely bound up with other subtle auxiliary investigations, which in turn lead to a remarkable extension of the theory of circle division. It is chiefly this wonderful interconnection of truths. that, as has often been observed, gives higher arithmetic such a distinctive appeal. Naturally these proofs themselves cannot be excerpted here; readers must consult the paper itself. But a couple of other new arithmetical theorems, also closely connected with the proof of the second criterion, deserve to be singled out here because of their simplicity.

\Needspace{24\baselineskip}Let $p$ be a prime of the form $4k+1$, and put this number $=ee+ff$
\SourceLeaf{178}%
so that $ee$ denotes the odd square and $ff$ the even square. If we further set
\[
1\pdot2\pdot3\mathinner{.\,.\,.\,.\,.\,.}\ k=q
\]
\[
(k+1)(k+2)(k+3)\dotsfive2k=r
\]
then $\pm e$ is always the least remainder obtained by dividing $\frac r{2q}$ by $p$, and $\pm f$ the least remainder obtained by dividing $\frac12rr$ by $p$ (where the least remainder is always understood to be taken between the bounds $-\frac12p$ and $+\frac12p$). The number $\frac r{2q}$, whose value for $p=5$ is $1$, can for larger values of $p$ also be put in the following form:
\[
\frac{6\pdot10\pdot14\pdot18\dotsfive(p-3)}{2\pdot3\pdot4\pdot5\dotsfive k}
\]
It is very remarkable that the decomposition of the number $p$ into two squares can thus be obtained quite directly. But an accompanying circumstance is almost more remarkable still. This procedure always gives the root of the odd square, $e$, with a positive sign if $e$, taken positively, has the form $4m+1$, and with a negative sign if $e$, taken positively, has the form $4m+3$. By contrast, no general rule whatever has yet been found for the sign with which the root of the even square, $f$, emerges from this operation, either a priori or by examining examples. At the end of the paper the author therefore recommends this subject to the friends of higher arithmetic for further investigation, convinced that success would also open up a fruitful source of new extensions of this beautiful part of mathematics.
\BookRule{.28\linewidth}
